% GENERATED FILE. Edit canonical lessons, then run npm run generate:books.
% canonical-source-hash: fnv1a64:601b8cc5014a62ea
% canonical-lessons: SA-C323-stupah, SA-C323-praciram, SA-C323-atitam, SA-C323-utkhanati, SA-C323-aviskaroti

\chapter{Ruins of the Past}
\label{ch:sa-ruins-of-the-past}
\hlchaptermodality{323}

\begin{chapteropening}
\textbf{By the end of this chapter:} \emph{I can describe old ruins and what was found there.}

Complete the last lesson of chapter 323: I can describe old ruins and what was found there.
\end{chapteropening}

\section[\texorpdfstring{\sk{स्तूपः} (stūpaḥ)}{स्तूपः (stūpaḥ)}]{\sk{स्तूपः} (stūpaḥ) — a stupa; a mound}

\label{lesson:SA-C323-stupah}



\begin{quote}
Before the new one: say the Sanskrit for founds, then the Sanskrit for introduces.
\end{quote}

\subsection*{You'll want to know: \sk{स्तूपः}}
\textbf{\sk{स्तूपः}} — \emph{stūpaḥ} — \textquotedblleft{}a stupa; a mound\textquotedblright{}.

\begin{quote}
\sk{पर्वतस्य} \sk{समीपे} \sk{प्राचीनः} \sk{स्तूपः} \sk{अस्ति।}
\end{quote}

\begin{quote}
\emph{parvatasya samīpe prācīnaḥ stūpaḥ asti.}
\end{quote}

\begin{quote}
There is an ancient stupa near the hill.
\end{quote}

A masculine noun.

\begin{cousinweb}
\textbf{\sk{स्तूपः}} (\emph{stūpaḥ}) MW gives \emph{stūpa} as a knot of hair, a top, and a heap of earth or bricks, especially a Buddhist monument, and reports two old derivations, so its origin is uncertain (MW).
\end{cousinweb}

\begin{grammarlens}[title={What you've built}]
One more word for telling what happened.
\end{grammarlens}

\subsection*{Your turn}
\begin{itemize}
\raggedright
  \item \emph{Say it:} \emph{stūpaḥ}
  \item \emph{Say it:} \emph{stūpaḥ}, once more
  \item \emph{Say it:} \emph{parvatasya samīpe prācīnaḥ stūpaḥ asti.} — There is an ancient stupa near the hill.
  \item \emph{Recall:} say the Sanskrit for founds, then the Sanskrit for introduces, then say \emph{stūpaḥ} again
\end{itemize}

\begin{tcolorbox}[breakable,colback=teal!4,colframe=teal!35!black,title={Before you move on}]
What does \sk{स्तूपः} mean? (\textquotedblleft{}A stupa; a mound\textquotedblright{}.) Say it once more.
\end{tcolorbox}
\section[\texorpdfstring{\sk{प्राचीरम्} (prācīram)}{प्राचीरम् (prācīram)}]{\sk{प्राचीरम्} (prācīram) — a rampart; an enclosing wall}

\label{lesson:SA-C323-praciram}



\begin{quote}
Before the new one: say the Sanskrit for introduces, then the Sanskrit for a stupa.
\end{quote}

\subsection*{You'll want to know: \sk{प्राचीरम्}}
\textbf{\sk{प्राचीरम्}} — \emph{prācīram} — \textquotedblleft{}a rampart; an enclosing wall\textquotedblright{}.

\begin{quote}
\sk{दुर्गस्य} \sk{प्राचीरं} \sk{जीर्णम्} \sk{अस्ति।}
\end{quote}

\begin{quote}
\emph{durgasya prācīraṃ jīrṇam asti.}
\end{quote}

\begin{quote}
The fort's rampart is crumbling.
\end{quote}

A neuter noun.

\begin{cousinweb}
\textbf{\sk{प्राचीरम्}} (\emph{prācīram}) is built on \emph{pra-cīra}, by MW's account; MW gives \emph{prācīra} as an enclosure, hedge, fence, wall (MW).
\end{cousinweb}

\begin{grammarlens}[title={What you've built}]
One more word for telling what happened.
\end{grammarlens}

\subsection*{Your turn}
\begin{itemize}
\raggedright
  \item \emph{Say it:} \emph{prācīram}
  \item \emph{Say it:} \emph{prācīram}, once more
  \item \emph{Say it:} \emph{durgasya prācīraṃ jīrṇam asti.} — The fort's rampart is crumbling.
  \item \emph{Recall:} say the Sanskrit for introduces, then the Sanskrit for a stupa, then say \emph{prācīram} again
\end{itemize}

\begin{tcolorbox}[breakable,colback=teal!4,colframe=teal!35!black,title={Before you move on}]
What does \sk{प्राचीरम्} mean? (\textquotedblleft{}A rampart; an enclosing wall\textquotedblright{}.) Say it once more.
\end{tcolorbox}
\section[\texorpdfstring{\sk{अतीतम्} (atītam)}{अतीतम् (atītam)}]{\sk{अतीतम्} (atītam) — the past}

\label{lesson:SA-C323-atitam}



\begin{quote}
Before the new one: say the Sanskrit for a stupa, then the Sanskrit for a rampart.
\end{quote}

\subsection*{You'll want to know: \sk{अतीतम्}}
\textbf{\sk{अतीतम्}} — \emph{atītam} — \textquotedblleft{}the past\textquotedblright{}.

\begin{quote}
\sk{इतिहासः} \sk{अतीतं} \sk{वर्णयति।}
\end{quote}

\begin{quote}
\emph{itihāsaḥ atītaṃ varṇayati.}
\end{quote}

\begin{quote}
History describes the past.
\end{quote}

A neuter noun.

\begin{cousinweb}
\textbf{\sk{अतीतम्}} (\emph{atītam}) is built from \emph{ati}, \textquotedblleft{}beyond\textquotedblright{}, and \emph{ita}, \textquotedblleft{}gone\textquotedblright{}, from the root √i, \textquotedblleft{}to go\textquotedblright{}; MW gives \emph{atīta} as gone by, past, and as a noun the past (MW).
\end{cousinweb}

\begin{grammarlens}[title={What you've built}]
One more word for telling what happened.
\end{grammarlens}

\subsection*{Your turn}
\begin{itemize}
\raggedright
  \item \emph{Say it:} \emph{atītam}
  \item \emph{Say it:} \emph{atītam}, once more
  \item \emph{Say it:} \emph{itihāsaḥ atītaṃ varṇayati.} — History describes the past.
  \item \emph{Recall:} say the Sanskrit for a stupa, then the Sanskrit for a rampart, then say \emph{atītam} again
\end{itemize}

\begin{tcolorbox}[breakable,colback=teal!4,colframe=teal!35!black,title={Before you move on}]
What does \sk{अतीतम्} mean? (\textquotedblleft{}The past\textquotedblright{}.) Say it once more.
\end{tcolorbox}
\section[\texorpdfstring{\sk{उत्खनति} (utkhanati)}{उत्खनति (utkhanati)}]{\sk{उत्खनति} (utkhanati) — digs up; excavates}

\label{lesson:SA-C323-utkhanati}



\begin{quote}
Before the new one: say the Sanskrit for a rampart, then the Sanskrit for the past.
\end{quote}

\subsection*{You'll want to know: \sk{उत्खनति}}
\textbf{\sk{उत्खनति}} — \emph{utkhanati} — \textquotedblleft{}digs up; excavates\textquotedblright{}.

\begin{quote}
\sk{वैज्ञानिकाः} \sk{भूमिम्} \sk{उत्खनन्ति।}
\end{quote}

\begin{quote}
\emph{vaijñānikāḥ bhūmim utkhananti.}
\end{quote}

\begin{quote}
The scientists dig up the ground.
\end{quote}

Past (he …): \sk{सः} \sk{भूमिम्} \sk{उत्खातवान्}, \emph{saḥ bhūmim utkhātavān}.

\begin{cousinweb}
\textbf{\sk{उत्खनति}} (\emph{utkhanati}) is built from \emph{ud}, \textquotedblleft{}up, out\textquotedblright{}, and the root √khan, \textquotedblleft{}to dig\textquotedblright{}, as in \emph{khanati}; MW gives \emph{ut-khan} as to dig up or out, excavate (MW).
\end{cousinweb}

\begin{grammarlens}[title={What you've built}]
One more word for telling what happened.
\end{grammarlens}

\subsection*{Your turn}
\begin{itemize}
\raggedright
  \item \emph{Say it:} \emph{utkhanati}
  \item \emph{Say it:} \emph{utkhanati}, once more
  \item \emph{Say it:} \emph{vaijñānikāḥ bhūmim utkhananti.} — The scientists dig up the ground.
  \item \emph{Recall:} say the Sanskrit for a rampart, then the Sanskrit for the past, then say \emph{utkhanati} again
\end{itemize}

\begin{tcolorbox}[breakable,colback=teal!4,colframe=teal!35!black,title={Before you move on}]
What does \sk{उत्खनति} mean? (\textquotedblleft{}Digs up; excavates\textquotedblright{}.) Say it once more.
\end{tcolorbox}
\section[\texorpdfstring{\sk{आविष्करोति} (āviṣkaroti)}{आविष्करोति (āviṣkaroti)}]{\sk{आविष्करोति} (āviṣkaroti) — brings to light; discovers}

\label{lesson:SA-C323-aviskaroti}



\begin{quote}
Before the new one: say the Sanskrit for the past, then the Sanskrit for digs up.
\end{quote}

\subsection*{You'll want to know: \sk{आविष्करोति}}
\textbf{\sk{आविष्करोति}} — \emph{āviṣkaroti} — \textquotedblleft{}brings to light; discovers\textquotedblright{}.

\begin{quote}
\sk{वैज्ञानिकः} \sk{प्राचीनं} \sk{नगरम्} \sk{आविष्करोति।}
\end{quote}

\begin{quote}
\emph{vaijñānikaḥ prācīnaṃ nagaram āviṣkaroti.}
\end{quote}

\begin{quote}
The scientist brings an ancient city to light.
\end{quote}

Past (he …): \sk{सः} \sk{नगरम्} \sk{आविष्कृतवान्}, \emph{saḥ nagaram āviṣkṛtavān}.

\begin{cousinweb}
\textbf{\sk{आविष्करोति}} (\emph{āviṣkaroti}) is built from \emph{āvis}, \textquotedblleft{}openly, before the eyes\textquotedblright{}, and the root √kṛ, \textquotedblleft{}to make\textquotedblright{}; MW gives \emph{āviṣ-kṛ} as to make apparent, reveal, uncover, and to discover is its modern use (MW).
\end{cousinweb}

\begin{grammarlens}[title={What you've built}]
One more word for telling what happened.
\end{grammarlens}

\subsection*{Your turn}
\begin{itemize}
\raggedright
  \item \emph{Say it:} \emph{āviṣkaroti}
  \item \emph{Say it:} \emph{āviṣkaroti}, once more
  \item \emph{Say it:} \emph{vaijñānikaḥ prācīnaṃ nagaram āviṣkaroti.} — The scientist brings an ancient city to light.
  \item \emph{Recall:} say the Sanskrit for the past, then the Sanskrit for digs up, then say \emph{āviṣkaroti} again
\end{itemize}

\begin{tcolorbox}[breakable,colback=teal!4,colframe=teal!35!black,title={Before you move on}]
What does \sk{आविष्करोति} mean? (\textquotedblleft{}Brings to light; discovers\textquotedblright{}.) Say it once more.
\end{tcolorbox}
